\lecture{17}{Thu 05 Dec 2024 09:30}{Group Actions}
For example, consider $D_n$ acting on $\left\{ 1,\ldots, n \right\} $ by arranging $\left\{ 1, \ldots,n \right\} $ as the vertices of a regular $n$-gon. This is a group action. We can also consider a group $G$ acting on itself by left multiplication; $\Pi_g : G \to G$ defined as $a\mapsto ga$. Notice this is basically just Cayley's theorem. One more example is that the set $S_n$ acts on the set $X\coloneqq \left\{ 1,\ldots,n \right\} $ by $\Pi_\sigma (i)=\sigma(i)$ for all $i\in X$.
\begin{proposition}
	Let $G$ be a group that acts on a set $X$. Then the group of permutations on the set $X$ given by the group action is a group.
\end{proposition}
\begin{proof}
	Multiplication follows equivalently to the group $G$ because of the secnond property of group actions, and the identity permutation is guarenteed to be defined properly for inverses to exist, so this proof is very trivial.
\end{proof}
\begin{notation}
	When our reference to $\Pi_g$ is clear, we write $\Pi_g(x) $ as $g\cdot x$, or in some places, $gx$.
\end{notation}
\begin{notation}
	If $G$ acts on $X$, we write $G \circlearrowright X$ or $G\curvearrowright X$.
\end{notation}
\begin{theorem}\label{thm:smnd}
	Let $G\circlearrowright X$. If $x\in X$, $g\in G$, and $y=g\cdot x$, then $x=g^{-1}\cdot y$. If $x\neq X'$, then $g\cdot x\neq g\cdot x'$.
\end{theorem}
\begin{proof}
	$g^{-1}\cdot y=g^{-1}\left( g\cdot x \right) =g^{-1}g\cdot x=e\cdot x=x$. Also, if $g\cdot x=g\cdot x'$, then $g^{-1}g\cdot x=g^{-1}g\cdot x'$, and thus $e\cdot x=e\cdot x'$, a contradiction.
\end{proof}
\begin{notation}
	$\operatorname{Sym}(X) $ is defined as the set of permutations of the set $X$. It is just a representation of $S_{\left| X \right| }$ given the precise elements of $X$.
\end{notation}
\begin{theorem}\label{thm:kjA}
	Actions of a group $G$ on a set $X$ are the same as group homomorphisms from $G$ to $\operatorname{Sym}(X) $.
	\[
		G\circlearrowright X\rightsquigarrow \phi : G \to \operatorname{Sym}(X) 
	.\]
\end{theorem}
\begin{proof}
	no proof!
\end{proof}

For example, consider $\operatorname{GL}_2(\mathbb{R}) \circlearrowright \mathbb{R}^2$. The action is basically for any  $A\in \operatorname{GL}_2(\mathbb{R}) $ and $\mathbf{x}\in \mathbb{R}$, we have $\Pi_A(\mathbf{x})=A \mathbf{x}$. Since $\operatorname{GL}_2(\mathbb{R}) =\operatorname{Iso}(\mathbb{R}^2) $, this holds. These sets are infinite and it still behaves very nicely. As another example, $S_n\circlearrowright R[x_1, \ldots, x_n]$ for $R$ a ring. For $\sigma \in S_n$, the action is
\[
	\sigma \cdot f(x_1, \ldots,x_n)=f(x_{\sigma (1)}, \ldots, x_{\sigma (n)})
.\]
The map $\sigma $ just acts on the indicies of the variables. A specific  example of this example is the following: consider $S_3$ and $\mathbb{Z}[x_1,x_2,x_3]$. Let $(12)\in S_3$. We have
\[
	(12)\cdot \left( x_1^2+x_2+x_3 \right) =x_2^2+x_1+x_3
.\]
The polynomials fixed by this action are called \textbf{symmetric} \textbf{polynomials}.
\begin{definition}[Symmetric Polynomials]\label{dfn:SPDN}
	Let $S_n\circlearrowright R[x_1, \ldots, x_n]$. The set of all $p(x_1, \ldots, x_n)$ such that for all $\sigma \in S_n$, the polynomial is fixed $\sigma \cdot p(x_1, \ldots, x_n)=p(x_1, \ldots, x_n)$, is called the set of \textbf{symmetric} \textbf{polynomials}.
\end{definition}
For example, the set of symmetric polynomials in $\mathbb{Z}[x_1,x_2,x_3]$ look like the set generated by
\[\left\{ x_1x_2x_3,x_1x_2+x_1x_3+x_2x_3,x_1+x_2+x_3 \right\} .\]
\begin{definition}\label{dfn:kjsd}
	Let $G\circlearrowright X$. For each $x\in X$, the orbit  of $x$ is
	\[
		\operatorname{Orb}_x = \left\{ g\cdot x \mid g\in G \right\}\subseteq X  
	.\]
	The stabilizer of $x$ is the set
	\[
		\operatorname{Stab}_x=\left\{ g\in G \mid g\cdot x=x \right\} \subseteq G 
	.\]
	We call $x$ a \textbf{fixed} \textbf{point} if $\operatorname{Stab}_x=G $.
\end{definition}
Note that the set of symmetric polynomials in $R[x_1, \ldots, x_n]$ is just the set of all fixed points in $R[x_1, \ldots, x_n]$.

Consider $\operatorname{GL}_2(\mathbb{R}) \circlearrowright \mathbb{R}^2$. The orbit
\[
	\operatorname{Orb}_{(0,0)}=\left\{ M\cdot \begin{pmatrix}
		0\\
		0
	\end{pmatrix}  \mid M\in \operatorname{GL}_2(\mathbb{R})  \right\} =\left\{ \begin{pmatrix}
		0\\
		0
	\end{pmatrix}  \right\} 
.\]
The stabilizer
\[
	\operatorname{Stab}_{(0,0)}=\left\{ M\in \operatorname{GL}_2(\mathbb{R})  \mid M\cdot \begin{pmatrix}
		0\\
		0
	\end{pmatrix} =\begin{pmatrix}
		0\\
		0
	\end{pmatrix}  \right\} =\operatorname{GL}_2(\mathbb{R})  
.\]
For $G\circlearrowright G$, for all $x\in G$, $\operatorname{Orb}_x=G $, and $\operatorname{Stab}_x=\left\{ e \right\}  $.

Consider now $G\circlearrowright G$ by congujation; that is, $g\cdot a=gag^{-1}$. The orbit of any element $a$ is the set
\[
	\operatorname{Orb}_a=\left\{ gag^{-1} \mid g\in G \right\} =cl(a)
.\]
The stabilizer of $a$ is the set
\[
	\operatorname{Stab}_a=\left\{ g\in G \mid gag^{-1}=a \right\} =C(a) 
.\]
Finally,  $a$ is a fixed point if and only if $C(a)=G$, and thus $a\in Z(G)$.
\begin{theorem}[Fundamental Theorem of Group Actions]\label{thm:the theorem}
	Let $G$ be a group, $X$ be a set, and let $G\circlearrowright X$. Then the following are true.
	\begin{enumerate}
		\item Different orbits are disjoint, so orbits partition $X$.
		\item For all $x\in X$, the stabilizer $\operatorname{Stab}_x $ is a subgroup of $G$. Also, $\operatorname{Stab}_{gx}=g \operatorname{Stab}_x g^{-1} $.
		\item For all $x\in X$, there exists a bijection
			\[
				G /\operatorname{Stab}_x \to \operatorname{Orb}_x
			.\]
		In particular, $\left| G \right| / \operatorname{Stab}_x = \left| \operatorname{Orb}_x \right| $. Equivalently, $\left| G \right| =\left| \operatorname{Orb}_x \right| \left| \operatorname{Stab}_x \right| $, the orbit stabilizer theorem.
	\end{enumerate}
\end{theorem}
